Nos últimos anos, o uso de técnicas de machine learning (ML) e visão computacional na agricultura tem se expandido significativamente, oferecendo soluções avançadas para monitoramento, classificação e colheita de produtos cultivados em ambientes controlados. Em especial para este estudo, a aplicação dessas técnicas no cultivo de cogumelos tem se mostrado promissora. 

Diversos estudos têm explorado algoritmos de detecção de objetos para localizar cogumelos em imagens e determinar o estágio de crescimento, visando otimizar a colheita.\cite{MOYSIADIS2023} utilizaram o algoritmo YOLOv5 para classificar cogumelos em três estágios de crescimento, chegando a obter uma acurácia de 70\% para o estágio final, quando os cogumelos estão prontos para a colheita. O estudo demonstrou que a classificação do sistema pode contribuir para decisões mais eficientes sobre o momento ideal, evitando perdas por amadurecimento excessivo ou coleta precoce. 

Outros estudos relevantes incluem a utilização do SSD (Single Shot MultiBox Detector) em robôs colhedores de cogumelos, que alcançaram até 95\% de acurácia na detecção e 86,8\% de sucesso na colheita individual, evidenciando a viabilidade prática de sistemas automatizados em ambientes controlados. Pesquisas voltadas para a classificação de espécies de cogumelos comestíveis e venenosos mostraram que ensembles de classificadores (Decision Tree, KNN, Random Forest, SVM, Naive Bayes) podem atingir 100\% de acurácia, reforçando a importância da inteligência artificial na diferenciação de cogumelos comestíveis e tóxicos. 

A integração de Internet das Coisas (IoT) com ML em sistemas de fungicultura inteligente também tem recebido destaque. Sensores de temperatura e umidade, associados a microcontroladores, permitem ajustar automaticamente variáveis críticas para o crescimento do micélio e do corpo de frutificação. Algoritmos de ML garantem a identificação precisa de cogumelos comestíveis, prevenindo a produção e consumo de exemplares tóxicos. Resultados experimentais indicam que esses sistemas mantêm condições ideais de temperatura e umidade, reduzem o esforço manual e fornecem dados em tempo real para os produtores. Avaliações de usabilidade (SUS) também demonstraram alto nível de satisfação, com pontuação média de 82\%, evidenciando a aplicabilidade prática dessas soluções. 

De forma complementar, \cite{Rahman2021} demonstrou que a integração de sensores ambientais com sistemas IoT permite monitoramento contínuo em tempo real e emissão de alertas automáticos, garantindo respostas rápidas a alterações ambientais e promovendo uma gestão mais eficiente do processo produtivo.  

Além do monitoramento, a classificação de espécies de cogumelos comestíveis e venenosos é um foco relevante na literatura. \cite{MetlekCetiner2023} desenvolveram um estudo utilizando um conjunto de dados com 8.124 amostras e 22 atributos visuais e físicos. Foram aplicados os algoritmos Random Forest (RF), Decision Tree (DT) e Regressão Logística (LR), com otimização de hiperparâmetros via GridSearchCV. Os resultados demonstraram que o RF apresentou melhor desempenho, alcançando precisão, recall e F1-score de até 0,97, 0,98 e 0,95, respectivamente, enquanto a DT otimizada destacou-se na distinção da classe venenosa. A regressão logística apresentou desempenho competitivo, mas ligeiramente inferior. O estudo evidencia que a seleção adequada de atributos e a otimização de parâmetros são fundamentais para reduzir erros na identificação de espécies. 

Dessa forma, a integração entre algoritmos de aprendizagem profunda para análise visual, sensoriamento ambiental e tecnologias de Internet das Coisas (IoT) para monitoramento em tempo real configura uma abordagem tecnológica robusta e de elevada relevância para a produção de cogumelos. Soluções convencionais frequentemente apresentam restrições de acesso para usuários e produtores em campo, decorrentes de limitações de infraestrutura e da complexidade operacional. O sistema proposto neste projeto trata tais barreiras, viabilizando monitoramento remoto contínuo e fornecendo suporte automatizado à tomada de decisão. Essa abordagem estabelece uma base sólida para o desenvolvimento de sistemas inteligentes aplicáveis diretamente ao cultivo de espécies como o shimeji, promovendo incremento na produtividade, melhoria na segurança alimentar e maior sustentabilidade da fungicultura.